% !PRINTED 307
\weberpage{307}

For the first assumption $(a, b, c) = (1, 0, 0)$ the determinant
of $\sigma$ becomes equal to $b_1 c_2 - c_1 b_2$, which can become $= 1$ in six
different ways, namely:

\[
\begin{pmatrix} b_1, & c_1 \\ b_2, & c_2 \end{pmatrix} = \begin{pmatrix} 1, & 0 \\ 0, & 1 \end{pmatrix}, \begin{pmatrix} 1, & 0 \\ 1, & 1 \end{pmatrix}, \begin{pmatrix} 0, & 1 \\ 1, & 0 \end{pmatrix}, \begin{pmatrix} 0, & 1 \\ 1, & 1 \end{pmatrix}, \begin{pmatrix} 1, & 1 \\ 1, & 0 \end{pmatrix}, \begin{pmatrix} 1, & 1 \\ 0, & 1 \end{pmatrix}.
\]

Then $a_1, a_2$ can further be assumed in four different
ways. Hence, when the first row is $(1, 0, 0)$,
one has 24 different forms of $\sigma$. The same number results
for the assumption of the first row in the forms $(0, 1, 0), (0, 0, 1)$.
But one obtains the same for the other assumptions
about the first row. For if this is, e.\,g., $(0, 1, 1)$, then

\[
a_1 (b_2 - c_2) - a_2 (b_1 - c_1) \equiv 1 \text{ (mod 2)}
\]

must hold, which again yields six possibilities for $a_1, b_1 - c_1, a_2, b_2 - c_2$,
and since one can assume $c_1, c_2$ in four ways, one
again obtains 24 possibilities. Finally, when the
first row is $(1, 1, 1)$,

\[
(a_1 - c_1)(b_2 - c_2) - (a_2 - c_2)(b_1 - c_1) \equiv 1 \text{ (mod 2)}
\]

must hold, which leads to the same number. The total number of all different
substitutions $\sigma$, i.\,e. the degree of $S$, is therefore 168.
The number 168 as the degree number of a group has already
met us once, namely with the group $L_7$ (\S.\,66, 72), and hence
the conjecture suggests itself that the groups $S$ and $L_7$
might be isomorphic. If this conjecture is confirmed,
the investigation of the group $S$ would thereby be essentially
facilitated, since we already know the divisors of the
group $L_7$.

But it is Theorem I. in \S.\,72 that serves to test this
conjecture.

To this end we seek to select generating elements $\chi, \omega, \Theta, \tau$
of the group $S$ in such a way that they satisfy the characteristic
conditions of Theorem I., \S.\,72. This can be done in
several ways. A general procedure for it is, to be sure, lacking,
but it succeeds easily by a few trials.

One first chooses an element of degree 7 arbitrarily
and forms from it the period, say
%CONT>>
